\scanpage{10}
124\hfill J.-P. \textsc{Serre}

\textbf{Remarques.} 1) Le groupe $\Phi$ est un $C$-sous-groupe algébrique connexe de $\mathbf{GL}_V(C)$. Il en résulte que $M$ est connexe, donc \textit{contenu dans la composante neutre} de $G_{\mathrm{alg}}$.

2) Soit $d\varphi$ l'endomorphisme de $V_C$ qui est l'homothétie de rapport $i$ sur $V_C(i)$. Les théorèmes 1 et 2, traduits en termes d'algèbres de Lie, signifient que:
\[
d\varphi\in\mathfrak g_{\mathrm{alg}}\otimes C.
\]

3) Soit $K_{nr}$ l'extension non ramifiée maximale de $K$ contenue dans $\overline K$. L'image $I$ de $\operatorname{Gal}(\overline K/K_{nr})$ dans $G$ est le \textit{sous-groupe d'inertie} de $G$. Comme $I$ est distingué dans $G$, $V$ est un $I$-module semi-simple et le groupe $I_{\mathrm{alg}}$ est défini. \textit{On a $M\subset I_{\mathrm{alg}}$}; cela se voit en appliquant le th. 2 sur le corps de base $\widehat K_{nr}$ complété de $K_{nr}$. (On n'aurait donc rien perdu si l'on avait supposé le corps résiduel $k$ algébriquement clos.)
\margin{En fait, on a $M=I_{\mathrm{alg}}$ (Sen)}

\begin{center}\textbf{\S\ 4. Un cas particulier}\end{center}

Conservons les notations des \S\S\ 2, 3, et faisons sur $V$ les hypothèses suivantes:

(H*.1) \textit{$V$ est un $\mathfrak g$-module absolument simple.}

(H*.2) \textit{$V_C$ est somme directe de $V_C(0)$ et $V_C(1)$.}

(H*.3) \textit{Les dimensions $n_0$ et $n_1$ de $V_C(0)$ et $V_C(1)$ sont $\geq1$ et premières entre elles.}

\textbf{Remarques.} -- 1) Les hypothèses (H*.1) et (H*.2) entraînent évidemment les hypothèses (H.1) et (H.2) du \S\ 3.

2) (H*.1) équivaut à dire que le $\mathfrak g$-module $V$ est \textit{semi-simple} et que \textit{son commutant est réduit aux homothéties} (\textsc{Bourbaki}, \textit{Alg.} VIII, \S\ 13, n\textsuperscript{o} 4, cor. à la prop. 5). En particulier, si $\mathfrak c$ désigne le centre de $\mathfrak g$, on voit que $\mathfrak c$ est, ou bien réduit à 0, ou bien égal à l'ensemble des homothéties de $V$. Dans les deux cas, c'est une algèbre de Lie algébrique, et l'on a donc $\mathfrak g_{\mathrm{alg}}=\mathfrak g$.

\textbf{Théorème 3.} -- \textit{Sous les hypothèses ci-dessus, on a $G_{\mathrm{alg}}=\mathbf{GL}_V$.}

Vu la remarque 2) ci-dessus, cet énoncé équivaut à:

\textbf{Corollaire 1.} -- \textit{L'algèbre de Lie $\mathfrak g$ de $G$ est égale à $\operatorname{End}(V)$.}

Il équivaut aussi à:

\textbf{Corollaire 2.} -- \textit{Le groupe $G$ est un sous-groupe ouvert de $\operatorname{Aut}(V)$.}

(En d'autres termes, l'action de $\operatorname{Gal}_K$ sur $V$ est aussi peu triviale que possible.)

\textbf{Démonstration.} -- Posons, pour simplifier $E=V_C$, $E_0=V_C(0)$ et
\margin{9}
